\chapter{Information-Theoretic Perspective}
\label{chap:information-theory}

% ── Emergent gravity ────────────────────────────────────────────────────────

\section{Emergent Gravity from Information Theory}
\label{sec:emergent-gravity-info}

\begin{itemize}
  \item Link to Verlinde's entropic gravity programme and related work on
        emergent gravity from an information-theoretic perspective.
  \item Connection to the entropic-gravity signatures discussed in
        Section~\ref{sec:entropic-gravity}.
\end{itemize}

% ── Local Planck length ─────────────────────────────────────────────────────

\section{Local Planck Length as Minimal Information Unit}
\label{sec:planck-info-unit}

\begin{itemize}
  \item Local Planck length $\ell_P(p)$ as the minimal quantisable size of an
        information unit at a point $p \in M$.
  \item Relation to the substrate parameter: $\ell_P \sim \hbar_q$.
\end{itemize}

% ── Time and information transfer ───────────────────────────────────────────

\section{Time as Information Transfer}
\label{sec:time-info}

\begin{itemize}
  \item Time as a measure of information transfer and change.
  \item Time slows at heat death ($q \to 1$) because information transfer
        rate decreases.
  \item Time runs fast at the Big Bang ($q \to 0$, equivalently $\hbar_q \to 1$) because information
        transfer is high; torsion of $E$ provides the driving mechanism.
\end{itemize}

% ── Big Bang perspective ────────────────────────────────────────────────────

\section{Big Bang as Initial Information State}
\label{sec:big-bang-info}

\begin{itemize}
  \item At $t = 0$: no information is known about the state of the QOU;
        matter is concentrated with high torsion pressure.
  \item At $t > 0$: the bundle $E$ acts on fermions, information begins to
        propagate, and the brane expands.
\end{itemize}

% ── Planck units and computational mechanics ────────────────────────────────

\section{Planck Units and Computational Mechanics}
\label{sec:planck-computation}

% TODO: develop the analogue of the Planck length for time and clarify its
% role in computational mechanics.

